\section{Rich tubes lie in a short list of plates}
\label{sec:geometry}

We next isolate a geometric consequence of two simultaneous properties:
every tube sees substantial mass, and the tubes available at each point
have a spread of directions. Both properties must concern the same
measure. The conclusion is a short list of planes near which all the
tube trajectories lie. The first part of the proof finds a rich plane
through each target line; the second part packs these planes, keeping
track of their offsets as well as their normals.

Divide all three physical coordinates by $L$, put $H=L$, and divide
the original measure by $mL^2$. In these coordinates the ball hypothesis
becomes $\nu(B(q,r))\lesssim r^2$ for $r\geq1$, the relevant domain
has diameter $O(H)$, and the tubes have unit width. This section uses
these coordinates throughout.

\begin{proposition}[Rich-tube geometry]\label{prop:rich}
Fix bounded velocity, domain, ball, and neighborhood constants,
a richness constant $c_{\mathrm r}>0$, an angular constant
$C_{\mathrm a}>0$, exponents $A_{\mathrm r},A_{\mathrm a}\geq0$,
and $\eta>0$.
There are constants $A_*,C_*,C'$ depending only on these fixed data
such that the following holds whenever $X\geq2$ and
$H\geq C_*X^{A_*}$.

Let $\nu$ be a nonnegative measure supported in a ball of radius $O(H)$,
with
\[
 \nu(B(q,r))\leq C_{\mathrm b}r^2\qquad(r\geq1).
\]
Let $\mathcal T$ be a finite family of lines
\[
 \ell_T=\{(b_T+2\omega_Tt,t):t\in\R\},\qquad
 |\omega_T|\leq C_{\mathrm v},
\]
viewed as unit-width tubes over a fixed time interval of length $O(H)$
containing the time projection of $\supp\nu$. Write
$d_T(x,t)=|x-b_T-2\omega_Tt|$ and
$e_T=(2\omega_T,1)/|(2\omega_T,1)|$. Assume:
\begin{enumerate}[label=\textup{(\roman*)}]
\item Every tube is rich against $\nu$:
\begin{equation}\label{eq:geometry-rich-hypothesis}
 \int (1+d_T(q))^{-M}\,d\nu(q)
 \geq c_{\mathrm r} H X^{-A_{\mathrm r}},\qquad M=100.
\end{equation}
\item At $\nu$-almost every $q$ there is a measurably chosen probability
$(p_T(q))_{T\in\mathcal T}$, supported on tubes with
$d_T(q)\leq C_{\mathrm n}$, such that every spherical direction ball
satisfies
\begin{equation}\label{eq:geometry-angular-hypothesis}
 \sum_{e_T\in B_{\mathbb S^2}(e,r)}p_T(q)
 \leq C_{\mathrm a}X^{A_{\mathrm a}}r^\eta
 \qquad(r\geq H^{-1}).
\end{equation}
\end{enumerate}
Then each tube can be assigned to one of at most $X^{C'}$ plates
(planes thickened by $X^{C'}$). Its assigned plate contains its
central line segment over the relevant time interval. If $n$ is a
unit normal to that plate, then
\begin{equation}\label{eq:src11}
 |n\cdot(2\omega_T,1)|\leq X^{C'}/H.
\end{equation}
The constants are independent of the number of tubes. In particular,
the richness in \eqref{eq:geometry-rich-hypothesis} and the pointwise
law in \eqref{eq:geometry-angular-hypothesis} use exactly the same
measure $\nu$.
\end{proposition}

Fixed factors may be absorbed into fixed powers of $X$, since $X\geq2$.
We write $X^{O(1)}$ only for powers depending on the data in the
proposition; the order in which all additional exponents are chosen
is specified at the end of the proof. No estimate below uses a ball
bound at radii smaller than one.
Bounded velocities also give
$\dist(q,\ell_T)\leq d_T(q)\leq C\dist(q,\ell_T)$, so transverse
deviation and orthogonal distance define comparable neighborhoods.
For a finite family, measurability of the neighboring law adds no
restriction: the set $\{T:d_T(q)\leq C_{\mathrm n}\}$ takes only
finitely many values, and one may choose an admissible probability
once for each occurring incidence pattern.
These patterns are Borel, since each $d_T$ is continuous, so the
resulting law is measurable.

\begin{lemma}[A truncated parameter energy]\label{lem:geometry-energy}
Let $H\geq1$ and let $\sigma$ be a finite measure on $\R^2$ of mass
$m$, satisfying
\[
 \sigma(B(z,r))\leq A r^{1+\eta}
 \qquad(z\in\R^2,\ r\geq H^{-1}),\qquad A>0.
\]
Set $\gamma=\eta/(1+\eta)$. Then
\begin{equation}\label{eq:geometry-parameter-energy}
 \iint \min\{H,|z-z'|^{-1}\}\,d\sigma(z')\,d\sigma(z)
 \leq C_\eta A^{1/(1+\eta)}m^{1+\gamma},
\end{equation}
where the kernel on the diagonal is $H$.
\end{lemma}

\begin{proof}
The assertion is immediate when $m=0$. Otherwise put
$R=(m/A)^{1/(1+\eta)}$. If $R<H^{-1}$, the inner integral is at most
\[
 Hm\leq m/R=A^{1/(1+\eta)}m^\gamma.
\]
If $R\geq H^{-1}$, the contribution of $|z-z'|\leq H^{-1}$ is at
most $AH^{-\eta}$. The dyadic annuli between $H^{-1}$ and $R$
contribute at most
\[
 C\sum_{H^{-1}\leq 2^jH^{-1}\lesssim R}
       A(2^jH^{-1})^\eta
 \leq C_\eta AR^\eta.
\]
The remaining region contributes at most $m/R$. Thus the inner
integral is bounded by
$C_\eta A^{1/(1+\eta)}m^\gamma$ in both cases. Integrating in $z$
proves \eqref{eq:geometry-parameter-energy}. The cutoff at $H$ is
what permits the first case without a hypothesis below $H^{-1}$.
\end{proof}

\begin{proof}[Proof of Proposition~\ref{prop:rich}]
\emph{1. Localize richness and choose a longitudinal point.}
For $1\leq w\leq H$, a width-$w$ neighborhood of any relevant
trajectory can be covered by $O(H/w)$ balls of radius $O(w)$.
Consequently
\begin{equation}\label{eq:geometry-tube-cover}
 \nu\{d_T\leq w\}\lesssim Hw.
\end{equation}
For $w\geq H$ the same upper bound follows from
$\nu(\R^3)\lesssim H^2$. Therefore the part of the smoothed
integral outside width $W\geq1$ is at most
\[
 C\sum_{j\geq0}(2^jW)^{-M}H(2^{j+1}W)
 \lesssim HW^{1-M}.
\]
Choose $W=X^{c_1}$, with $(M-1)c_1>A_{\mathrm r}$ by a sufficiently
large fixed margin. It follows from
\eqref{eq:geometry-rich-hypothesis} that, for every $T$,
\begin{equation}\label{eq:geometry-local-richness}
 \nu\{d_T\leq W\}\geq cH X^{-A_{\mathrm r}}.
\end{equation}
We always take $H$ large enough that $W\ll H$.

Fix a target tube $T_0$, and abbreviate its line and unit direction
by $\ell$ and $e$. Translate the origin to a point of $\ell$ at
distance $O(W)$ from the original domain. Such a point exists by
\eqref{eq:geometry-local-richness}. The measure domain, and all
relevant trajectory segments, remain inside a ball of radius $O(H)$:
the latter assertion follows for every tube from its own localized
richness, bounded velocity, and the common $O(H)$ time range.

Draw $q_0$ according to normalized $\nu$ restricted to
$\{d_{T_0}\leq W\}$, and let
$s=H^{-1}e\cdot q_0$. Thus $Hs$ is the longitudinal coordinate of
the orthogonal projection of $q_0$ onto $\ell$. A longitudinal window
of relative length $r\geq H^{-1}$ inside this width-$O(W)$ tube
is covered by $O(1+Hr/W)$ balls of radius $O(W)$. Before normalization
its mass is at most $C(W^2+HrW)$; hence
\begin{equation}\label{eq:geometry-longitudinal}
 \Prob\{s\in I\}\lesssim
 X^{A_{\mathrm r}}(W^2/H+rW)
 \leq X^{O(1)}r,
 \qquad |I|\leq r,\quad r\geq H^{-1}.
\end{equation}
In the last inequality we used precisely $H^{-1}\leq r$.

\emph{2. Sample a transverse neighbor and keep bin measures
unnormalized.}
Choose $\delta=X^{-c_0}$, with $c_0$ large relative to
$A_{\mathrm a}/\eta$. At every sampled point discard directions
with $|e\times e_T|<\delta$. They lie in the union of two spherical
balls about $e$ and $-e$ of radius $O(\delta)$, so
\eqref{eq:geometry-angular-hypothesis} makes their total probability
less than $1/2$. Here $H^{-1}\leq\delta$ is included in the scale
hierarchy. Renormalize the remaining probabilities separately at
each $q_0$. This costs at most two in the angular bound and does not
change the marginal law of $q_0$ or \eqref{eq:geometry-longitudinal}.
Draw a tube $T$ from this remaining law. It is a physical neighbor
of $q_0$, and is rich against the original common measure $\nu$.

Partition azimuth about the target axis $e$ into half-open intervals
of length comparable to $H^{-1}$. For a bin $\pi$, let $u_\pi$ be
its central radial unit vector in $e^\perp$, and let $P_\pi$ be the
exact plane through $\ell$ spanned by $e,u_\pi$. For a sampled unit
direction $e'$ in this bin define
\[
 a=\frac{e\cdot e'}{u_\pi\cdot e'}.
\]
The denominator is at least $c\delta$, so $|a|\lesssim\delta^{-1}$.
Let $\sigma_\pi$ be the pushforward to $(s,a)$ of the sampling
probability \emph{restricted} to this azimuth bin, and write
$m_\pi=\|\sigma_\pi\|$. We have
\begin{equation}\label{eq:geometry-bin-masses}
 \sum_\pi m_\pi=1.
\end{equation}
In particular, we have not divided $\sigma_\pi$ by $m_\pi$.

For completeness, the angular estimate remains effective in these
coordinates without conditioning on the bin. If $n_\pi$ completes
$e,u_\pi$ to an orthonormal basis and $\theta$ is azimuth relative
to $u_\pi$, then
\[
 e'=\frac{ae+u_\pi+\tan\theta\,n_\pi}
             {(1+a^2+\tan^2\theta)^{1/2}},
 \qquad |\theta|\lesssim H^{-1}.
\]
An interval of length $O(r)$ in $a$, within this bin, is therefore
contained in a direction ball of radius $O(r)$ when $r\geq H^{-1}$;
allowing $X^{O(1)}r$ would also suffice. For a ball in the
$(s,a)$-plane of radius $r$, integrate the conditional angular bound
over the corresponding longitudinal event in
\eqref{eq:geometry-longitudinal}. This gives
\begin{equation}\label{eq:geometry-joint-frostman}
 \sigma_\pi(B(z,r))\leq X^{A_{\mathrm F}}r^{1+\eta}
 \qquad(r\geq H^{-1})
\end{equation}
for any $A_{\mathrm F}>A_{\mathrm r}+2c_1+A_{\mathrm a}$ with a
sufficiently large fixed margin. This is a bound for the
unnormalized measure. No assertion about angular nonconcentration
conditional on landing in a small bin is needed.

\emph{3. Strips in the bin plane.}
Write $x=e\cdot q$ and $y=u_\pi\cdot q$. Choose
$w=X^{c_2}$ with $c_2>c_0+c_1$ by a sufficiently large fixed margin,
and define the plate and its strip multiplicity by
\begin{align}
 Q_\pi&=\{q\in\supp\nu:\dist(q,P_\pi)\leq w\},\notag\\
 F_\pi(q)&=\int
  \mathbf1_{\{|x-Hs-ay|\leq w\}}\,d\sigma_\pi(s,a).
 \label{eq:geometry-strip-multiplicity}
\end{align}
For each sampled pair $(q_0,T)$ in bin $\pi$, the whole width-$W$
rich portion of $T$ lies in $Q_\pi$ and in its associated strip.
Indeed, write
$q=q_0+\tau e_T+\varepsilon$ with $|\varepsilon|\lesssim W$ and
$|\tau|\lesssim H$. Azimuth binning gives
$|n_\pi\cdot e_T|\lesssim H^{-1}$, while
$e\cdot e_T-a u_\pi\cdot e_T=0$. Since the displacement of $q_0$
from its projection $Hs e$ is $O(W)$, the plate error is $O(W+1)$
and the strip error is $O((1+|a|)W)$, both bounded by $w$.

Consider two strips with parameters $z=(s,a)$ and $z'=(s',a')$.
At an intersection point,
\begin{equation}\label{eq:geometry-strip-intersection}
 |H(s-s')+(a-a')y|\leq2w,\qquad |y|\lesssim H.
\end{equation}
In particular,
\[
 |s-s'|\leq C|a-a'|+2w/H.
\]
Put $d=|z-z'|$. If $d\leq Cw/H$, covering the full strip by unit
balls costs $X^{O(1)}H$, which is at most
$X^{O(1)}\min\{H,d^{-1}\}$ after increasing the fixed power.
If $d>Cw/H$, with $C$ chosen from the domain bound, the existence
of an intersection in \eqref{eq:geometry-strip-intersection}
forces $|a-a'|\gtrsim d$. The allowed $y$-interval then has length
$O(w/d)$. Its cross-section, the plate thickness, and the factor
$1+|a|$ cost only fixed powers of $X$ in a unit-ball covering.
The extra constant in the count of unit intervals is also absorbed,
because $d\leq X^{O(1)}$. If the offset difference dominates so
strongly that \eqref{eq:geometry-strip-intersection} cannot hold,
the intersection is empty. Thus in all cases the intersection has
a covering by
\begin{equation}\label{eq:geometry-strip-kernel}
 X^{A_{\mathrm I}}\min\{H,|z-z'|^{-1}\}
 \quad\hbox{unit balls}
\end{equation}
For example, the covering argument gives
$Cw^3(1+\delta^{-1})\min\{H,d^{-1}\}$, so any
$A_{\mathrm I}>3c_2+c_0$ with a sufficiently large fixed margin
suffices in \eqref{eq:geometry-strip-kernel}.

Use the ball bound at radius one, Fubini's theorem, and
Lemma~\ref{lem:geometry-energy} with
\eqref{eq:geometry-joint-frostman}. For
$\gamma=\eta/(1+\eta)$ this yields
\begin{equation}\label{eq:geometry-strip-energy}
 \int_{Q_\pi}F_\pi^2\,d\nu
 \leq X^{A_{\mathrm E}}m_\pi^{1+\gamma}.
\end{equation}
One can take any sufficiently enlarged exponent
$A_{\mathrm E}\geq A_{\mathrm I}+A_{\mathrm F}/(1+\eta)$.
There is no remaining factor of $H$ in this estimate.

\emph{4. Richness away from the target and overlap of the far plates.}
Put $R_{\mathrm f}=H/X^{c_3}$, choosing $c_3$ after $c_0,c_1,c_2$,
and let
\[
 Q_\pi^{\mathrm{far}}
   =Q_\pi\cap\{q:\dist(q,\ell)\geq R_{\mathrm f}\}.
\]
The geometric meaning of the strip and this exclusion is shown in
Figure~\ref{fig:geometry-far}.
\begin{figure}[tbp]
\centering
\begin{tikzpicture}[scale=0.9,>=Stealth,font=\small]
 \fill[black!4] (-0.5,-0.7) rectangle (6.2,0.7);
 \draw[densely dashed,black!45] (-0.5,0.7)--(6.2,0.7);
 \draw[densely dashed,black!45] (-0.5,-0.7)--(6.2,-0.7);
 \draw[->] (-0.5,0)--(6.5,0) node[right] {$e$};
 \draw[->] (0,-1.25)--(0,2.85) node[above] {$u_\pi$};
 \node[below] at (5.15,0) {target line $\ell$};
 \fill[blue!9] (0.5,-1.1)--(0.95,-1.1)--(4.7,2.5)--(4.25,2.5)--cycle;
 \draw[blue!65!black,thick,->] (0.73,-1.1)--(4.48,2.5)
       node[above right] {sampled direction $e'$};
 \draw[blue!50!black,densely dashed] (0.5,-1.1)--(4.25,2.5);
 \draw[blue!50!black,densely dashed] (0.95,-1.1)--(4.7,2.5);
 \fill (1.88,0) circle (1.5pt);
 \node[below right] at (1.88,0) {$Hs$};
 \fill[blue!65!black] (3.55,1.61) circle (2pt);
 \node[right] at (3.67,1.5) {retained far mass};
 \draw[<->,black!65] (5.8,0.04)--(5.8,0.68)
       node[midway,right] {$R_{\mathrm f}$};
 \node[anchor=west] at (0.15,2.65) {$P_\pi$};
 \node[anchor=west,blue!65!black] at (0.4,-1.6)
       {strip: $|x-Hs-ay|\leq w$};
\end{tikzpicture}
\caption{Geometry in the bin plane $P_\pi$ (schematic projection).
The shaded horizontal band is the near region removed in
\eqref{eq:geometry-near-mass}; the thin sloping strip contains the
sampled tube after enlargement. Since its direction is separated
from the target, only a controlled length lies in the near band.
The true plate also has thickness $w$ perpendicular to this picture.}
\label{fig:geometry-far}
\end{figure}

Because a sampled tube makes sine-angle at least $\delta$ with
$\ell$, the length of its width-$W$ portion within distance
$R_{\mathrm f}$ of $\ell$ is at most
$C\delta^{-1}(W+R_{\mathrm f})$. Covering it by radius-$O(W)$
balls gives the mass bound
\begin{equation}\label{eq:geometry-near-mass}
 C\delta^{-1}W(W+R_{\mathrm f})
 \leq C X^{c_0+2c_1}
       +CHX^{c_0+c_1-c_3}.
\end{equation}
Choose $c_3>A_{\mathrm r}+c_0+c_1$ by a large fixed margin, and
then require $H\gg X^{A_{\mathrm r}+c_0+2c_1}$. The bound in
\eqref{eq:geometry-near-mass} is less than half the lower bound in
\eqref{eq:geometry-local-richness}, uniformly for every sampled
tube. Integrating their retained richness against the sampling
law therefore gives
\begin{equation}\label{eq:geometry-strip-first-moment}
 \int_{Q_\pi^{\mathrm{far}}}F_\pi\,d\nu
 \geq cHX^{-A_{\mathrm r}}m_\pi.
\end{equation}
This step uses the richness of every neighboring tube against
$\nu$, including mass far from the point at which it was sampled.

Cauchy--Schwarz, \eqref{eq:geometry-strip-energy}, and
\eqref{eq:geometry-strip-first-moment} imply
\begin{equation}\label{eq:geometry-far-plate-mass}
 \nu(Q_\pi^{\mathrm{far}})
 \geq H^2X^{-K_0}m_\pi^{1-\gamma}
\end{equation}
for some $K_0\geq2A_{\mathrm r}+A_{\mathrm E}$, enlarged to absorb
fixed constants.

All planes $P_\pi$ pass through the same line $\ell$. At a point
$q$ at distance $\rho\geq R_{\mathrm f}$ from that line, the
condition $\dist(q,P_\pi)\leq w$ restricts $u_\pi$ to two azimuth
arcs of length $O(w/\rho)$. The two arcs account for opposite radial
directions. Since bins have size comparable to $H^{-1}$, their
overlap multiplicity is at most
\[
 C(1+Hw/R_{\mathrm f})\leq X^{A_{\mathrm O}},
 \qquad A_{\mathrm O}>c_2+c_3
\]
after a fixed enlargement. Take $H\gg X^{c_2+c_3}$ so that
$w\ll R_{\mathrm f}$. Summing
\eqref{eq:geometry-far-plate-mass}, using this overlap and
$\nu(\R^3)\lesssim H^2$, gives
\begin{equation}\label{eq:geometry-bin-concentration}
 \sum_\pi m_\pi^{1-\gamma}\leq X^{K_1}
\end{equation}
for a fixed $K_1>K_0+A_{\mathrm O}$. But
\[
 1=\sum_\pi m_\pi
 \leq (\max_\pi m_\pi)^\gamma
          \sum_\pi m_\pi^{1-\gamma}.
\]
Hence some bin has $m_\pi\geq X^{-K_1/\gamma}$. Its plane contains
the exact target line, and its width-$w$ plate has mass at least
\begin{equation}\label{eq:geometry-rich-plane}
 H^2X^{-K},\qquad
 K\geq K_0+(1-\gamma)K_1/\gamma.
\end{equation}
We have produced such a rich plane for every target tube.

\emph{5. Pack rich planes using both normal and offset separation.}
Return to a common origin at the center of the original domain.
Write each plane just obtained as
\[
 P_j=\{q:n_j\cdot q=b_j\},\qquad |n_j|=1,
\]
with $(n_j,b_j)$ identified with $(-n_j,-b_j)$. The width-$w$
plates, restricted to $\supp\nu$, all have mass at least
$\lambda H^2$, where $\lambda=X^{-K}$. Let
$C_{\mathrm t}$ be a fixed bound for $\nu(\R^3)/H^2$.

Choose an angular threshold $\vartheta=X^{c_4}/H$, with
$c_4>2c_2+2K$ by a sufficiently large fixed margin. If the
projective angle between two normals is at least $\vartheta$,
the intersection of their width-$w$ plates inside the domain is
contained in a box with dimensions
\[
 O(H)\ \times\ O(w)\ \times\
 O\bigl(w/|\sin\angle(n_i,n_j)|\bigr).
\]
Here projective angles lie in $[0,\pi/2]$, and $\vartheta\ll1$.
Unit-ball covering thus gives
\begin{equation}\label{eq:geometry-plate-intersection}
 \nu(Q_i\cap Q_j)
 \leq C H^2X^{2c_2-c_4}
 \leq \varepsilon H^2,
 \qquad \varepsilon\leq\lambda^2/(2C_{\mathrm t}).
\end{equation}
For a pair of normals at projective angle below $\vartheta$, align
their signs. Changing the normal then changes its defining linear
function by at most $C H\vartheta$ on the domain. Choose an offset
threshold
\begin{equation}\label{eq:geometry-offset-threshold}
 \Delta=X^{c_5}>2w+CH\vartheta,
 \qquad c_5>\max\{c_2,c_4\}
\end{equation}
with sufficiently large fixed margins. Such a close-normal pair
with $|b_i-b_j|\geq\Delta$ has disjoint plates on the domain.

Select a maximal sublist in which every pair is separated either
by projective normal angle at least $\vartheta$, or, after aligning
closer normals, by offset at least $\Delta$. Finiteness makes this
selection immediate. If the list has length $N$, set
$G=\sum_{j=1}^N\mathbf1_{Q_j}$. Its first and second moments obey
\[
 \int G\,d\nu\geq N\lambda H^2,
 \qquad
 \int G^2\,d\nu\leq C_{\mathrm t}NH^2+\varepsilon N^2H^2.
\]
Cauchy--Schwarz and the total mass bound give
\[
 N^2\lambda^2
 \leq C_{\mathrm t}^2N+C_{\mathrm t}\varepsilon N^2.
\]
By the choice in \eqref{eq:geometry-plate-intersection}, the last
term is absorbed, and
\begin{equation}\label{eq:geometry-plate-count}
 N\leq2C_{\mathrm t}^2\lambda^{-2}\leq X^{O(1)}.
\end{equation}

Maximality says that each original plane is close to a selected
one in both its projective normal and its aligned offset. On the
entire bounded range of relevant trajectory segments, this means
that its central line lies within
$\Delta+CH\vartheta=X^{O(1)}$ of the selected plane. Enlarge the
selected plates accordingly. If $n_0$ is the original normal of
the plane through the target line and $n$ is the aligned selected
normal, then
\[
 n_0\cdot(2\omega_T,1)=0,\qquad |n-n_0|\lesssim\vartheta.
\]
Bounded velocity now gives
$|n\cdot(2\omega_T,1)|\lesssim X^{c_4}/H$, proving
\eqref{eq:src11}. Offset control is essential for containment of
the segment; angular control supplies the asserted directional
bound.

\emph{Constant hierarchy.}
Here is one complete order of choices. First fix the data in the
proposition and $M=100$. Choose $c_0,c_1$ so that
\[
 c_0\eta>A_{\mathrm a},\qquad
 (M-1)c_1>A_{\mathrm r},
\]
with fixed margins that give the two required half-mass losses.
Next choose $c_2>c_0+c_1$, followed by
\[
 A_{\mathrm F}>A_{\mathrm r}+2c_1+A_{\mathrm a},\qquad
 A_{\mathrm I}>3c_2+c_0,\qquad
 A_{\mathrm E}>A_{\mathrm I}+A_{\mathrm F}/(1+\eta),
\]
with sufficient fixed margins for
\eqref{eq:geometry-joint-frostman}--\eqref{eq:geometry-strip-energy}.
Choose
\[
 c_3>\max\{c_2,A_{\mathrm r}+c_0+c_1\},
\]
and then choose $A_{\mathrm O},K_0,K_1,K$ as in
\eqref{eq:geometry-far-plate-mass}--\eqref{eq:geometry-rich-plane}.
Choose $c_4>2c_2+2K$ to ensure
\eqref{eq:geometry-plate-intersection}, and then
$c_5>\max\{c_2,c_4\}$ to ensure
\eqref{eq:geometry-offset-threshold}. Choose $C'$ larger than
$c_5,c_4,2K$, with a fixed margin absorbing the final widths,
counts, and direction constants. Finally take
\[
 A_*>\max\{A_{\mathrm r}+c_0+2c_1,\ c_2+c_3,\ c_0,\ c_4,\ C'\},
\]
and $C_*$ sufficiently large. Every strict inequality here includes
an additional fixed margin if necessary for the fixed constants.
All choices precede $X,H$; in particular no exponent depends on
the number of tubes or on an azimuth-bin mass.
\end{proof}
